% language: swedish
\begin{example}[3.4]
Beräkna enhetsnormal till ytan $f(x, y, z) = x^2 + y^2 - z = 0$ i punkt $(1, 1, 2)$
\begin{align*}
\nabla f &= \nabla(x^2 + y^2 - z) = (2x, 2y, -1) \\
\nabla f(1, 1, 2) &= (2, 2, -1)
\end{align*}
Vektorn måste normeras $\abs{\nabla f(1, 1, 2)} = \sqrt{2^2 + 2^2 + 1} = 3$\\
Enhetsnormalen: $\displaystyle \hat{\vb{n}}\ \frac{\nabla f(1, 1, 2)}{\abs{\nabla f(1, 1, 2)}} = \left( \frac{2}{3}, \frac{2}{3}, -\frac{1}{3} \right)$
\end{example}

\subsubsection*{Gradient, konservativa fält och potentialer}
\begin{theorem}[3.1]
Relaterar gradient av skalärt fält till konservativa vektorfält.

Antag vektorfält $\vb{F}$ relaterar till skalärfält $\phi$ som $\vb{F} = \nabla\phi$.\\
$\nabla\phi$ existerar överallt i $D$.\\
Då är $\vb{F}$ konservativt i $D$ och omvänt om $\vb{F}$ är konservativt kan man skriva $\vb{F} = \nabla\phi$. $\phi$ kallas potential.
\end{theorem}

\begin{proof}
Antag $\vb{F} = \nabla\phi$. Beräkna linjeintegralen från $A$ till $B$ längs $C$.
\[ \int_C \vb{F} \cdot d\vb{r} = \int_C \nabla\phi \cdot d\vb{r} = \int_C d\phi = \Big[ \phi \Big]_A^B = \phi(B) - \phi(A) \]
$\Rightarrow \vb{F}$ är konservativt.

Omvänt: Antag $\vb{F}$ är konservativt
\[ \phi(\vb{r}) = \int_0^{\vb{r}} \vb{F} \cdot d\tilde{\vb{r}}, \quad \text{Vi vet att } d\phi = \vb{F} \cdot d\vb{r} \]
\begin{align*}
&\Rightarrow \vb{F} \cdot d\vb{r} = \nabla\phi \cdot d\vb{r} \text{ måste gälla } \forall d\vb{r} \\
&\Rightarrow \vb{F} = \nabla\phi
\end{align*}
\end{proof}

\begin{note}
Notera $\phi$ är ej unik, en konstant kan adderas.
\end{note}
